\section{Introduction}\label{sec:intro}
\pagenumbering{arabic} \setcounter{page}{1}

Singapore's power grid serves roughly 5.9 million residents in a tropical climate where air-conditioning runs year-round; national electricity consumption exceeded 55\,TWh in 2024 \parencite{ema2025statistics}. Because cooling is the single largest end-use in buildings, even small shifts in temperature or humidity can noticeably swing demand \parencite{zhao2012review}, and holidays, weekends, and monsoon transitions add further irregularity. For grid operators, getting short-term forecasts right is not optional---it directly affects generation scheduling and system reliability.

Classical approaches like ARIMA and exponential smoothing assume linear, stationary dynamics and tend to fall short once non-linear weather--demand interactions come into play \parencite{almusaylh2018forecasting}. Tree-based machine learning models handle these interactions much better, but their performance depends heavily on hyperparameter choices, and manual tuning is both tedious and often sub-optimal. Zhang et al.\ \parencite{zhang2023catboostppso} demonstrated that pairing CatBoost with Phasor Particle Swarm Optimisation (PPSO) beats five other meta-heuristics---however, their experiment used hourly data from a \emph{single residential customer} and only tuned the favoured model. Two questions naturally follow: does the approach still work at the scale of \emph{daily national-level} demand in a tropical climate, and does CatBoost-PPSO retain its edge when \emph{every} competing model is given the same tuning budget?

We address both questions. Specifically, we (i) construct a daily Singapore dataset by merging EMA demand records with Open-Meteo weather data (February 2023 -- December 2025); (ii) engineer 16 features spanning weather, derived-comfort, calendar, and autoregressive categories; (iii) compare eight model configurations---a Lag-1 baseline, Linear Regression, Random Forest, XGBoost, and CatBoost, plus PPSO-tuned variants of each tree model under an \emph{identical} budget---across six evaluation metrics; and (iv) deploy the best model through a web portal (FastAPI + React) aimed at non-technical users. Intra-day resolution, sector-level decomposition, and deep learning architectures are left for future work.
